\section{Analytic regularity and quadrature for the retained completion}
\label{sec:quadrature}
The approximation theorem leaves a finite bounded integral.
Controlling its evaluation requires more than the real size
estimate in \cref{subsec:rs-input}. We now specialize throughout
to \(1\le T\le2\), \(p=\lfloor L/2\rfloor\), \(K=2p\).
The statements below concern the finite retained expression, not
an analytic continuation of the real Riemann--Siegel error theorem.

\subsection{Remove apparent poles before estimating coefficients}
The exact coefficient identities used here are those recorded in
Arias de Reyna, Part II, (29), (62), (86)--(89)
\cite{ariasII}, attributed there to the published Part I
\cite{ariasI}. The underlying Part I proof has not been independently
rechecked in this project. Part II contains preliminary precision
arguments; its unproved technical inequalities are not imported.
We use the named identities and derive the following bounds directly.

In the same complex normalization as \cref{subsec:rs-input}, set
\begin{subequations}\label{eq:rs-coefficient-identity}
\begin{align}
 F_0(\chi)&=
 \frac{\e^{\pi\ii(\chi^2/2+3/8)}
                     -\ii\sqrt2\cos(\pi\chi/2)}{2\cos(\pi\chi)},\\
 C_k(\chi,1/2)&=\pi^{-2k}
  \sum_{j=0}^{\lfloor3k/2\rfloor}
       \left(\frac{\pi}{2\ii}\right)^j
                  d_j^{(k)}F_0^{(3k-2j)}(\chi).
\end{align}
\end{subequations}
At \(\chi=n+1/2\) the two numerator terms are both
\(\ii(-1)^{n(n+1)/2}\). Every zero of the denominator is simple
and removable, so \(F_0\) is entire.
This fact alone would not make direct numerical division near
those points stable.

\begin{lemma}[Weighted coefficient recurrence]\label{lem:weighted-rs}
Let \(e_j^{(k)}=(3k-2j)!\,d_j^{(k)}\), with \(e_0^{(0)}=1\)
and out-of-range entries zero. For \(m=3k-2j>0\),
\begin{equation}
 e_j^{(k)}
 =\frac{(m-1)(m-2)}4e_j^{(k-1)}-e_{j-2}^{(k-1)}.
 \label{eq:weighted-rs}
\end{equation}
The first term is structurally absent when \(m=1,2\).
For positive even \(k\), with \(j=3k/2\), the endpoint is
\begin{equation}
 e_j^{(k)}=-\sum_{h=1}^{j}\frac{(-1)^h}{h!}e_{j-h}^{(k)}.
 \label{eq:rs-endpoint}
\end{equation}
The values and their positive recurrence majorants satisfy
\begin{equation}
 \max_j\abs{e_j^{(k)}}\le E_k:=8^k(k!)^2.
 \label{eq:weighted-rs-bound}
\end{equation}
\end{lemma}
\begin{proof}
Multiply the exact coefficient recurrence at \(\sigma=1/2\)
by \(m!\). The previous-order derivative indices are \(m-3\)
and \(m+1\), which gives \eqref{eq:weighted-rs}.
The exact endpoint relation gives \eqref{eq:rs-endpoint}.
The sum of absolute interior coefficients is at most
\((9/4)k^2\). The endpoint contributes a factor at most
\(\sum_{h\ge1}1/h!<2\) times the largest interior majorant.
Thus a row costs at most \((9/2)k^2\) times the preceding row.
Induction with \(8k^2\) proves \eqref{eq:weighted-rs-bound},
also for intermediate positive sums.
\end{proof}

\begin{proposition}[Explicit complex coefficient bounds]
\label{prop:rs-complex}
Put
\[
 B_F=\tfrac12\{\e^{2\pi}+\sqrt2\cosh(\pi/2)\}.
\]
For \(\abs{\Rea\chi}\le3/2\), \(\abs{\Ima\chi}\le1/2\),
\begin{equation}
 \abs{C_k(\chi,1/2)}\le
 D_k:=\frac{B_F}{1-\pi/8}
                   (64/\pi^2)^k(k!)^2.
 \label{eq:complex-Ck}
\end{equation}
At real \(\chi\in[-1,1]\), for each fixed \(s\ge0\),
\begin{equation}
 \abs{\partial_\chi^sC_k(\chi,1/2)}
 \le B_F\pi^{-2k}E_k
       \sum_{j=0}^{\lfloor3k/2\rfloor}
       (\pi/2)^j\frac{(3k-2j+s)!}{(3k-2j)!}.
 \label{eq:real-Ck-derivative}
\end{equation}
\end{proposition}
\begin{proof}
On the rectangle \(\abs{\Rea\chi}\le2,\abs{\Ima\chi}\le1\),
the quotient is bounded by \(B_F\) on its boundary:
the denominator has modulus at least \(2\) on the vertical
sides and at least \(2\sinh\pi\) on the horizontal sides.
Maximum modulus applies to the entire continuation.
Radius-\(1/2\) Cauchy disks then give
\(\abs{F_0^{(m)}}\le B_Fm!2^m\) on the smaller rectangle.
Substitute \eqref{eq:weighted-rs-bound} into
\eqref{eq:rs-coefficient-identity} and sum
\(\sum_{j\ge0}(\pi/8)^j\), proving \eqref{eq:complex-Ck}.
At real \(\chi\in[-1,1]\), radius-one Cauchy disks prove
\eqref{eq:real-Ck-derivative}.
\end{proof}
In particular \(\log^+D_k=O((k+1)\log(k+2))\).
The complex normalization is essential; the real cosine quotient
\(\mathcal C\) in \eqref{eq:short} is not substituted for \(F_0\).

\subsection{Holomorphic extension of each finite branch}
Split \(I_X\) at its exact thresholds \(4\pi n^2\).
On a branch \(J\), fix \(m\) and use
\(\chi_m(u)=1-2(\sqrt{u/(4\pi)}-m)\).
If \(A_m(u)\) denotes the finite complex expression inside
the real part in \eqref{eq:real-completion}, then
\[
 A_m(u)+\overline{A_m(\bar u)}
\]
is holomorphic wherever the displayed finite factors are, and
agrees with the real initial correction on \(J\).
Add the two analytic transfer arms. This extends the retained
branch \(Q_m\) holomorphically; it makes no assertion about the
discarded remainder off the real axis.

Use the Bernstein ellipse of parameter \(\rho=2\) for each \(J\).
It lies within \(\abs{u-X}\le2V\), \(\abs{\Ima u}\le3V/4\),
and every point on it is within \(V\) of a real point of \(J\).
The gates
\[
 2V\le X/4,\qquad V/\sqrt{2\pi X}\le1/2
\]
ensure that the branch stays in the coefficient rectangle of
\cref{prop:rs-complex}. Write
\begin{subequations}\label{eq:ellipse-data}
\begin{align}
 a_*&=\sqrt{(X-2V)/(4\pi)},\\
 H_M&=\exp(2V\abs{\ell_+'(X)}+2HV^2/X),&
 H_U&=\exp(V(L+1)/2),\\
 \mathcal T_X&=\sum_{n\in\mathcal T}n^{-1/2}\e^{V\log n},&
 B_Q&=2H_M\left(\mathcal T_X+
                    H_Ua_*^{-1/2}\sum_{k=0}^KD_ka_*^{-k}\right).
\end{align}
\end{subequations}
Here \(\mathcal T\) is the union of actual transfer indices
over the branches; it contains at most one index.
The gamma bound follows from \cref{lem:coeff}.
For \(U\), integrate its logarithmic derivative
\(-\ii\Log(u/(4\pi))/4\), imposing the eventual gate
\(\sup_{\abs{q}\le2V}\abs{\Log((X+q)/(4\pi))}\le L+1\).
Thus \(\abs{Q_m(u)/g}\le B_Q\) on every branch ellipse.

Define nonnegative kernel polynomials
\[
 \mathsf K_a(v)=a!\sum_{j=0}^{\lfloor a/2\rfloor}
       \frac{(v/2)^{a-2j}(1/4)^j}{(a-2j)!j!},
 \qquad 0\le a\le3,
\]
and
\begin{equation}
 K_E=\frac{\e^{(Y+3V/4)^2/4}}{\sqrt{4\pi}}
       \left[\max_{0\le a\le3}\mathsf K_a(W+2V)
                       +\mathsf K_2(W+2V)\right],
 \qquad B_E=B_QK_E.
 \label{eq:ellipse-kernel}
\end{equation}
The sum of the maximum spatial channel and the time channel
of the normalized integrand is bounded by \(B_E\).
For the latter use \(\partial_Tk_T=\partial_z^2k_T\);
these are actual kernel derivatives.
The estimates retain the off-real kernel growth, giving
\(\log^+B_E=O(L^2)+O(K\log(K+2))\).

\subsection{Analytic Gauss error}
\begin{lemma}[A complex-valued Gauss bound]\label{lem:gauss}
If an analytic integrand is bounded by \(B\) on the
Bernstein \(\rho\)-ellipse of a real interval \(J\), its
\(n\)-node Gauss--Legendre rule has error at most
\[
 4\abs J B\,\frac{\rho^{-2n+1}}{\rho-1}.
\]
The statement also holds for a finite-dimensional vector integrand
in any norm.
\end{lemma}
\begin{proof}
Under the Joukowski map, Cauchy's coefficient estimate gives
Chebyshev coefficients bounded by \(2B\rho^{-j}\).
Truncation at degree \(2n-1\) has uniform remainder at most
\(2B\rho^{-2n+1}/(\rho-1)\).
Gauss weights are positive, sum to \(\abs J\), and the rule
is exact through degree \(2n-1\) \cite{dlmfgauss}.
Subtract the polynomial; the integral and the quadrature
each cost at most \(\abs J\) times the uniform remainder.
The same argument applies to the channel norm. It uses no
real mean-value formula for complex-valued data.
\end{proof}
Summing branch lengths, \(\sum_J\abs J=2V\), and setting
\(\rho=2\) gives
\begin{equation}
 E_{\mathrm{quad}}\le16VB_E4^{-n}.
 \label{eq:quadrature-bound}
\end{equation}
For any assigned center-normalized budget \(\varepsilon/20\),
it is sufficient to take
\begin{equation}
 n\ge\max\left\{1,\left\lceil
           \frac{\log(320VB_E/\varepsilon)}{\log4}\right\rceil\right\}.
 \label{eq:quadrature-order}
\end{equation}
When \(-\log\varepsilon=O(L^2)\) and \(K=O(L)\), this is
\(n=O(L^2)\) on each of at most two branches.
Node, weight, and input-data errors have not been included
in \eqref{eq:quadrature-bound}; they are paid next.
